% !PRINTED 1
\weberpage{1}
\begin{center}
{\large\gesperrt{EINLEITUNG.}}

\rule{0.18\linewidth}{0.4pt}
\end{center}

Wir setzen bei unseren Betrachtungen die \gesperrt{natürlichen
Zahlen} $1,2,3,\ldots$ und die Regeln, nach denen mit diesen
Zahlen gerechnet wird, als bekannt und gegeben voraus. Die
fundamentalen Rechenarten, die sogenannten \gesperrt{vier Species},
sind die \gesperrt{Addition}, die \gesperrt{Multiplication}, zu der als
Wiederholung das \gesperrt{Potenziren} gehört, die \gesperrt{Subtraction}
und die \gesperrt{Division}. Die beiden ersten heissen die \gesperrt{directen
Rechenoperationen}; sie sind dadurch ausgezeichnet, dass sie im
Reiche der natürlichen Zahlen unbegrenzt ausgeführt werden können.
Die erste der indirecten oder inversen Operationen, die Subtraction,
lässt sich nur dann ausführen, wenn der Minuend grösser ist als der
Subtrahend.

Die Aufgabe der Division kann man auf zwei Arten auffassen. Bei der
ersten elementaren Auffassung wird gefragt, wie oft der Divisor im
Dividenden enthalten ist. Eine Zahl ist nicht in einer kleineren
enthalten. Ist aber der Dividend gleich oder grösser als der Divisor,
so giebt die Beantwortung der Frage einen \gesperrt{Quotienten} und in
den meisten Fällen einen Rest, der kleiner als der Divisor ist. Wenn
kein Rest bleibt, so sagt man, die Division geht auf, oder der
\gesperrt{Dividend} ist durch den \gesperrt{Divisor} theilbar, oder der
Divisor ist ein \gesperrt{Factor} oder \gesperrt{Theiler} der Zahl, die
den Dividenden bildet.

Diese Aufgabe, die sich schon auf den ersten Stufen der Rechenkunst
einstellt, führt zu einer tiefer liegenden Unterscheidung der Zahlen,
die das Fundament aller Zahlentheorie ist.

Da ein Factor nie grösser sein kann als die Zahl, deren Factor sie
ist, so hat jede Zahl nur eine endliche Anzahl von Factoren. Jede Zahl
ist durch $1$ und durch sich selbst theilbar, und eine Zahl, die sonst
keinen Theiler hat, heisst eine \gesperrt{Prim-}
%CONT>>
